\section{Momentum tilt between equities and bonds}\label{sec:thesis-math}

Let \(P^{E}_t\) and \(P^{B}_t\) be the adjusted closes of SPY (equities) and TLT (bonds), and let
\(L\) be the lookback (126 sessions). The trailing momentum of each asset at the close of session
\(t\) is
\begin{equation}
  m^{a}_t=\frac{P^{a}_{t}}{P^{a}_{t-L}}-1,\qquad a\in\{E,B\}.
  \label{eq:tm-momentum}
\end{equation}
The leader is \(\ell_t=\arg\max_a m^a_t\). With base equity weight \(w_0=0.6\) and tilt
\(\tau=0.2\), the target weights are
\begin{equation}
  w^{\ell_t}_t=w^{\ell_t}_0+\tau,\qquad
  w^{a}_t=w^{a}_0-\tau\ \ (a\ne\ell_t),\qquad
  \sum_a w^a_t=1.
  \label{eq:tm-weights}
\end{equation}
These use only \(\mathcal F_{t}\) at the close of session \(t\) and are applied to returns from
\(t+1\), consistent with Equation~\eqref{eq:design-timing}. The excess return over the static
60/40 reference portfolio is linear in the tilt:
\begin{equation}
  r^{\mathrm{strat}}_{t+1}-r^{\mathrm{static}}_{t+1}
  =\tau\,s_t\,\bigl(r^{E}_{t+1}-r^{B}_{t+1}\bigr),\qquad
  s_t=\begin{cases}+1,& \ell_t=E,\\-1,& \ell_t=B.\end{cases}
  \label{eq:tm-excess}
\end{equation}
Because the cost is proportional to turnover, the net excess return per period has expectation
\begin{equation}
  \E\!\left[r^{\mathrm{strat}}_{t+1}-r^{\mathrm{static}}_{t+1}\right]
  =\tau\,\E\!\left[s_t\,(r^{E}_{t+1}-r^{B}_{t+1})\right]-c\,\E[\mathrm{TO}_t],
  \label{eq:tm-edge}
\end{equation}
where \(\mathrm{TO}_t\) is turnover. A tilt only trades when the leader flips, and each flip moves
\(2\tau\) of weight in total. Hypothesis H2 (edge survives costs) is the statement that
Equation~\eqref{eq:tm-edge} is positive at the assumed \(c\); H1 concerns the Sharpe ratio of the
same series and is evaluated empirically in Section~\ref{sec:results}.
